% F07：本书原创矢量示意；依据所在节的定义、证明或明确模型。
\begin{figure}[htbp]
\centering
\begin{tikzpicture}[x=1cm,y=1cm,>=stealth,every node/.style={font=\small},line width=.65pt]

\draw[->,gray](0,0)--(5.7,0)node[right]{\(x\)};
\foreach \x/\n in {3/1,1.5/2,.75/4}{\draw[->,very thick,AnalysisBlue](\x,0)--(\x,1.8);\node[above]at(\x,1.8){\(1/\n\)};}
\draw[->,orange!80!black](2.8,2.6)--(.3,2.6);
\node at(2.5,3.25){\(\delta_{1/n}\Rightarrow\delta_0\)};
\draw[->,gray](7,0)--(12.7,0)node[right]{\(x\)};
\fill[AnalysisBlue!10](7,0)rectangle(9.3,1.4);
\node at(8.15,.65){固定紧区间};
\foreach \x in {8,10,12}{\draw[->,very thick,AnalysisBlue](\x,0)--(\x,1.8);}
\draw[->,orange!80!black](9,2.6)--(12.5,2.6);
\node at(9.8,3.25){\(\delta_n\) 淡收敛到零};
\node at(9.8,-.65){总质量始终为 \(1\)};

\end{tikzpicture}
\caption[原子测度的汇聚与质量逃逸]{原子测度的汇聚与质量逃逸。每根箭头都代表单位质量。逃逸族的紧支集测试趋零，但常函数一的积分不变，因此它不弱收敛到零测度。}
\label{b1:fig:visual-f07}
\end{figure}
